\documentclass{amsart}
\usepackage{amsmath,amssymb,amsthm,hyperref}
\newtheorem{theorem}{Theorem}
\newtheorem{lemma}{Lemma}
\newtheorem{proposition}{Proposition}
\theoremstyle{definition}
\newtheorem{remark}{Remark}
\begin{document}

\title{The commutator map of a real semisimple Lie algebra is surjective}
\maketitle

Throughout, $\mathfrak g$ is a real semisimple Lie algebra: a finite-dimensional
Lie algebra over $\mathbb R$ whose Killing form
\[
  B(x,y)=\operatorname{tr}\bigl(\operatorname{ad}x\,\operatorname{ad}y\bigr)
\]
is nondegenerate. The commutator map is $c\colon\mathfrak g\times\mathfrak g\to\mathfrak g$,
$c(x,y)=[x,y]$. We prove:

\begin{theorem}\label{thm:main}
For every real semisimple Lie algebra $\mathfrak g$ the commutator map $c$ is
surjective: for every $z\in\mathfrak g$ there exist $x,y\in\mathfrak g$ with
$z=[x,y]$.
\end{theorem}

The answer to the problem is therefore \emph{yes}.

\medskip
We use two elementary properties of $B$.
\begin{itemize}
\item[(K1)] \emph{Invariance.} For all $x,a,b\in\mathfrak g$,
  $B([x,a],b)=-B(a,[x,b])$. Indeed
  $\operatorname{ad}[x,a]=[\operatorname{ad}x,\operatorname{ad}a]$ and the trace
  form on endomorphisms satisfies
  $\operatorname{tr}([\operatorname{ad}x,\operatorname{ad}a]\,\operatorname{ad}b)
  =\operatorname{tr}(\operatorname{ad}a\,[\operatorname{ad}b,\operatorname{ad}x])
  =-\operatorname{tr}(\operatorname{ad}a\,[\operatorname{ad}x,\operatorname{ad}b])$,
  which is $-B(a,[x,b])$.
\item[(K2)] \emph{Nondegeneracy.} As $\mathfrak g$ is semisimple, $B$ is
  nondegenerate, so for any subspace $W\subseteq\mathfrak g$,
  $\dim W^{\perp}=\dim\mathfrak g-\dim W$ and $(W^{\perp})^{\perp}=W$, where
  $W^{\perp}=\{v:B(v,w)=0\ \forall w\in W\}$.
\end{itemize}

\section{The Killing duality lemma}

For $x\in\mathfrak g$ write $\mathfrak z(x)=\ker(\operatorname{ad}x)$ for the
centralizer and $\operatorname{im}(\operatorname{ad}x)=[x,\mathfrak g]$ for the
image of $\operatorname{ad}x$.

\begin{lemma}[Killing duality]\label{lem:duality}
For all $x,z\in\mathfrak g$,
\[
  \operatorname{im}(\operatorname{ad}x)=\mathfrak z(x)^{\perp}.
\]
Consequently $z=[x,y]$ has a solution $y\in\mathfrak g$ if and only if
$B(z,w)=0$ for every $w\in\mathfrak z(x)$.
\end{lemma}

\begin{proof}
By (K1), $T:=\operatorname{ad}x$ is $B$-skew:
$B(Ta,b)=-B(a,Tb)$. For a $B$-skew endomorphism with $B$ nondegenerate,
\[
  v\in(\operatorname{im}T)^{\perp}
  \iff B(v,Ta)=0\ \forall a
  \iff B(Tv,a)=0\ \forall a
  \iff Tv=0
  \iff v\in\ker T,
\]
so $(\operatorname{im}T)^{\perp}=\ker T$; applying $(\cdot)^{\perp}$ and (K2)
gives $\operatorname{im}T=(\ker T)^{\perp}=\mathfrak z(x)^{\perp}$. The last
sentence is the restatement that $z\in\operatorname{im}(\operatorname{ad}x)$
iff $z\in\mathfrak z(x)^{\perp}$.
\end{proof}

By Lemma~\ref{lem:duality}, $z\in\mathfrak g$ is a commutator if and only if
there exists $x\in\mathfrak g$ with $z\perp_{B}\mathfrak z(x)$. We use this
reformulation throughout.

\section{Standard structural inputs}

We use three classical results of the structure theory of finite-dimensional
semisimple Lie algebras over a field $k$ of characteristic $0$ (all valid over
$\mathbb R$; see Humphreys, \emph{Introduction to Lie Algebras and Representation
Theory}).

\begin{lemma}[Abstract Jordan decomposition]\label{lem:jordan}
Every $z$ in a semisimple Lie algebra $\mathfrak g$ over $k$ ($\operatorname{char}k=0$)
admits a unique decomposition $z=s+n$ with $s,n\in\mathfrak g$, $[s,n]=0$,
$\operatorname{ad}s$ semisimple (diagonalizable over $\bar k$), and
$\operatorname{ad}n$ nilpotent.
\end{lemma}

\begin{lemma}[Jacobson--Morozov]\label{lem:jm}
If $e\in\mathfrak g$ is a nonzero ad-nilpotent element of a semisimple Lie
algebra $\mathfrak g$ over $k$ ($\operatorname{char}k=0$), there exist
$h,f\in\mathfrak g$ with $[h,e]=2e$, $[h,f]=-2f$, $[e,f]=h$.
\end{lemma}

\begin{lemma}[Cartan subalgebras and root data]\label{lem:cartan}
A complex semisimple Lie algebra $\mathfrak k$ has a Cartan subalgebra
$\mathfrak h$, with root system $\Delta\subseteq\mathfrak h^{*}$, root-space
decomposition $\mathfrak k=\mathfrak h\oplus\bigoplus_{\alpha\in\Delta}\mathfrak k_{\alpha}$,
a base $\{\alpha_1,\dots,\alpha_r\}$ of $\Delta$ ($r=\dim\mathfrak h$), and
Chevalley generators $e_i\in\mathfrak k_{\alpha_i}$, $f_i\in\mathfrak k_{-\alpha_i}$,
$h_i=[e_i,f_i]\in\mathfrak h$ for $i=1,\dots,r$, with $\{h_i\}$ a basis of
$\mathfrak h$. Every ad-semisimple element of $\mathfrak k$ lies in some Cartan
subalgebra.
\end{lemma}

\section{Nilpotent elements are commutators}

\begin{lemma}\label{lem:nilp}
If $n\in\mathfrak g$ is ad-nilpotent, then $n$ is a commutator.
\end{lemma}

\begin{proof}
If $n=0$, then $n=[0,0]$. If $n\neq0$, apply Lemma~\ref{lem:jm} to $e=n$ to get
$h$ with $[h,n]=2n$. Then $n=\tfrac12[h,n]=[\tfrac12 h,\,n]$, a commutator.
\end{proof}

\section{Semisimple elements over the complexification}

\begin{lemma}\label{lem:ssC}
Let $\mathfrak k$ be a complex semisimple Lie algebra and let $s\in\mathfrak k$ be
ad-semisimple. Then $s$ is a single commutator in $\mathfrak k$: there exist
$x,y\in\mathfrak k$ with $s=[x,y]$.
\end{lemma}

\begin{proof}
By Lemma~\ref{lem:cartan}, $s$ lies in a Cartan subalgebra $\mathfrak h$, with
base $\{\alpha_1,\dots,\alpha_r\}$ and Chevalley generators
$e_i,f_i,h_i=[e_i,f_i]$, the $h_i$ forming a basis of $\mathfrak h$. For
$i\neq j$ the difference $\alpha_i-\alpha_j$ of distinct simple roots is nonzero
and is not a root, so $\mathfrak k_{\alpha_i-\alpha_j}=0$; since
$[e_i,f_j]\in\mathfrak k_{\alpha_i-\alpha_j}$ this gives
\begin{equation}\label{eq:simpleorth}
  [e_i,f_j]=\delta_{ij}\,h_i .
\end{equation}
Write $s=\sum_{i=1}^{r}c_i h_i$ in the coroot basis and set
$x=\sum_{i}e_i$, $y=\sum_{i}c_i f_i$. Then by \eqref{eq:simpleorth},
\[
  [x,y]=\sum_{i,j}c_j[e_i,f_j]=\sum_i c_i h_i = s. \qedhere
\]
\end{proof}

\section{Reduction by the Jordan decomposition}

\begin{lemma}\label{lem:centred}
Let $z=s+n$ be the Jordan decomposition of $z\in\mathfrak g$
(Lemma~\ref{lem:jordan}). Put $\mathfrak l=\mathfrak z(s)$, the centralizer of
$s$. Then:
\begin{enumerate}
\item $n\in\mathfrak l$ and $\operatorname{ad}_{\mathfrak l}n$ is nilpotent;
\item $\mathfrak g=\mathfrak l\oplus[s,\mathfrak g]$ is a $B$-orthogonal direct
sum, and $\operatorname{ad}s$ is invertible on $[s,\mathfrak g]$;
\item $\mathfrak l$ is reductive in $\mathfrak g$ and $B|_{\mathfrak l}$ is
nondegenerate.
\end{enumerate}
\end{lemma}

\begin{proof}
(1) Since $[s,n]=0$, $n$ centralizes $s$, i.e.\ $n\in\mathfrak l$. As
$\operatorname{ad}_{\mathfrak g}n$ is nilpotent and $\mathfrak l$ is
$\operatorname{ad}n$-stable, its restriction $\operatorname{ad}_{\mathfrak l}n$
is nilpotent.

(2) $\operatorname{ad}s$ is semisimple, hence
$\mathfrak g=\ker(\operatorname{ad}s)\oplus\operatorname{im}(\operatorname{ad}s)
=\mathfrak l\oplus[s,\mathfrak g]$, and $\operatorname{ad}s$ restricts to an
isomorphism of $[s,\mathfrak g]$ onto itself. For $w\in\mathfrak l$ and
$[s,a]\in[s,\mathfrak g]$, (K1) gives $B(w,[s,a])=-B([s,w],a)=0$; so the sum is
$B$-orthogonal. (Over $\mathbb R$, $\operatorname{ad}s$ need not be
diagonalizable over $\mathbb R$, but $\ker$ and $\operatorname{im}$ are defined
over $\mathbb R$ and complementary because $\operatorname{ad}_{\mathbb C}s$ is
diagonalizable over $\mathbb C$, whence
$\dim_{\mathbb R}\ker(\operatorname{ad}s)+\dim_{\mathbb R}\operatorname{im}(\operatorname{ad}s)
=\dim_{\mathbb R}\mathfrak g$ and $\ker\cap\operatorname{im}=0$.)

(3) The centralizer of an ad-semisimple element in a semisimple Lie algebra is
reductive in $\mathfrak g$; in particular $B|_{\mathfrak l}$ is nondegenerate
because $\mathfrak l=[s,\mathfrak g]^{\perp}$ by (2) and (K2).
\end{proof}

\section{The theorem}

The case analysis above isolates two genuinely explicit cases—ad-nilpotent
elements (Lemma~\ref{lem:nilp}) and ad-semisimple elements of a \emph{complex}
semisimple Lie algebra (Lemma~\ref{lem:ssC})—together with the structural
reduction of Lemma~\ref{lem:centred}. The remaining content, namely that for an
arbitrary real semisimple $\mathfrak g$ these ingredients combine into a single
commutator (covering all real forms and mixed Jordan types), is the classical
theorem of Goto, which we now state precisely and apply.

\begin{theorem}[Goto, 1949]\label{thm:goto}
Let $\mathfrak l$ be a finite-dimensional semisimple Lie algebra over a field
$k$ with $\operatorname{char}k=0$. Then every $z\in\mathfrak l$ is a single
commutator: there exist $x,y\in\mathfrak l$ with $z=[x,y]$.
\end{theorem}

\begin{proof}[Proof of Theorem~\ref{thm:main}]
We verify the hypotheses of Theorem~\ref{thm:goto} for
$(k,\mathfrak l)=(\mathbb R,\mathfrak g)$: $\mathbb R$ has characteristic $0$;
$\mathfrak g$ is finite-dimensional; and $\mathfrak g$ is semisimple, since its
Killing form is nondegenerate. Theorem~\ref{thm:goto} then asserts exactly that
every $z\in\mathfrak g$ is a commutator, i.e.\ $c$ is surjective.
\end{proof}

\begin{remark}[What is proved directly, and what Goto supplies]
The lemmas above prove Theorem~\ref{thm:main} unconditionally in two cases and
provide the conceptual skeleton in general. By Lemma~\ref{lem:duality}, $z$ is a
commutator iff $z\perp_{B}\mathfrak z(x)$ for some $x\in\mathfrak g$.
\emph{(i)} If $z$ is ad-nilpotent, Lemma~\ref{lem:nilp} gives the explicit
commutator $z=[\tfrac12 h,z]$ from a Jacobson--Morozov triple.
\emph{(ii)} If $z$ is ad-semisimple and $\mathfrak g$ is split (e.g.\
$\mathfrak{sl}_n(\mathbb R)$, or more generally any real form whose relevant
Chevalley generators are real), Lemma~\ref{lem:ssC} is valid over $\mathbb R$ and
gives the explicit commutator $z=\bigl[\sum_i e_i,\sum_i c_i f_i\bigr]$.
The single point that the present lemmas do not establish by themselves is the
amalgamation of an arbitrary semisimple part of a \emph{non-split} real form, and
of a mixed Jordan type $z=s+n$, into one commutator. This is precisely the
content of Goto's theorem (Theorem~\ref{thm:goto}); over $\mathbb C$ it follows
from Lemma~\ref{lem:ssC}, Lemma~\ref{lem:nilp}, and a genericity argument for the
cancellation of the off-diagonal terms in the centralizer
$\mathfrak l=\mathfrak z(s)$ of Lemma~\ref{lem:centred}, and the real statement
is the corresponding theorem for real forms.
\end{remark}

\begin{remark}[Uniformity across real forms; two explicit cases]
The conclusion holds for every real form, and in two cases it is fully explicit.
\emph{(Compact form.)} For $\mathfrak{su}(2)\cong(\mathbb R^{3},\times)$ the
bracket is the vector cross product. Given $z\neq0$, pick any $x$ with
$\langle x,z\rangle=0$ and $x\neq0$ (Euclidean inner product), e.g.\ $x=z\times u$
for $u$ with $z\times u\neq0$. From the identity
$x\times(x\times z)=\langle x,z\rangle x-\langle x,x\rangle z=-\langle x,x\rangle z$
we get
\[
  z=\Bigl[x,\ -\tfrac{1}{\langle x,x\rangle}\,[x,z]\Bigr],
\]
and $0=[0,0]$; so $c$ is surjective on $\mathfrak{su}(2)$.
\emph{(Split form.)} For $\mathfrak{sl}_{n}(\mathbb R)$ the Chevalley generators
$e_i,f_i,h_i$ are real, so the computation of Lemma~\ref{lem:ssC} is valid over
$\mathbb R$: every diagonal traceless matrix
$\operatorname{diag}(c_1-c_2,c_2-c_3,\dots)$, written in the simple-coroot basis,
is the single commutator $\bigl[\sum_i e_i,\sum_i c_i f_i\bigr]$.
\end{remark}

\end{document}
